% !TeX root = ../../Book3.tex
\section{格与双周期函数}
\label{b3:sec:2201}

指数函数把一个平移周期变成单值函数，两个实线性无关的周期则把平面压成紧曲面。紧性迫使非常值函数具有极点，因此双周期理论自然属于亚纯函数，而非整函数。

\subsection{\texorpdfstring{\(\mathbb C\)}{C} 中的格}
\label{b3:subsec:220101}

\begin{definition}\label{b3:def:22-lattice}
若 \(\omega_1,\omega_2\in\mathbb C\) 在实数域上线性无关，则离散加法子群
\(\Lambda=\mathbb Z\omega_1+\mathbb Z\omega_2\) 称为复平面中的\term[ge]{格}[lattice]。有序生成元称为格的一组基；通常调整方向使 \(\operatorname{Im}(\omega_2/\omega_1)>0\)。
\end{definition}
实线性映射 \((x,y)\mapsto x\omega_1+y\omega_2\) 可逆，故存在 \(c,C>0\)，使
\(c\sqrt{m^2+n^2}\le|m\omega_1+n\omega_2|\le C\sqrt{m^2+n^2}\)。因此格点没有有限聚点，半径 \(R\) 内有 \(O(1+R^2)\) 个格点，且
\[
\sum_{0\ne\omega\in\Lambda}|\omega|^{-p}<\infty\quad(p>2).
\]
最后结论可按整数方框壳 \(\max(|m|,|n|)=j\) 求和；该壳有 \(8j\) 个点，故归结于 \(\sum j^{1-p}\)。这条估计将负责椭圆级数的正规收敛。

\subsection{基本平行四边形}
\label{b3:subsec:220102}

集合 \(P=\{s\omega_1+t\omega_2:0\le s,t<1\}\) 是\term[jiben-pingxing-sibianxing]{基本平行四边形}[fundamental parallelogram]。每个复数模 \(\Lambda\) 有唯一一个代表落在 \(P\)；闭平行四边形上对边点须按平移等同。

\ProseParagraphBreak

积分时常将整个平行四边形平移一个小量，使边界不经过给定亚纯函数的零极点。这样的平移确实存在：在固定有界范围内只有有限多个零极点，需要避开的平移参数是一族有限线段，其并集不含开集。计算最终只取决于模格的点，不取决于平行四边形的起点。

\subsection{双周期函数}
\label{b3:subsec:220103}

若亚纯函数 \(f\) 满足 \(f(z+\omega)=f(z)\) 对全部 \(\omega\in\Lambda\) 成立，就称它关于 \(\Lambda\) 双周期。检验两组生成元足够，因正向平移、反向平移和有限次复合给出全部格平移。

\ProseParagraphBreak

周期群可能比指定的格更大。例如相对于 \(\Lambda\) 的函数也可相对于子格 \(2\Lambda\) 讨论，但后一基本区域中极点数变成原来的四倍。凡说“次数”或“模格的零点”，必须先固定采用哪个格；不能无说明地用函数的全部周期替换给定格。

\subsection{椭圆函数}
\label{b3:subsec:220104}

上述亚纯双周期函数称为\term[tuoyuan-hanshu]{椭圆函数}[elliptic function]。常函数也包含在这个定义中。
\begin{proposition}\label{b3:prop:22-entire-constant}
没有非常值的整椭圆函数。
\end{proposition}
\begin{proof}
整函数在闭基本平行四边形上有界；周期性把这个界传到全平面。Liouville 定理使它为常数。
\end{proof}
仅有一个实周期时结论失败：\(e^{2\pi\mathrm{i}z}\) 是周期为一的非常值整函数，其商空间仍有无穷远端。两个独立周期使商空间紧致，才造成上述刚性。

\ProseParagraphBreak

离散性也能从周期本身反向检验。设非常值亚纯函数有非零周期 \(\omega_j\to0\)。在任一没有极点的小圆盘内，差商
\(\bigl(f(z+\omega_j)-f(z)\bigr)/\omega_j\) 恒为零，极限给 \(f'(z)=0\)。因此函数在一个开集上常值，亚纯恒等定理使它处处常值，矛盾。所以非常值亚纯函数的周期群一定离散。特别地，三个实线性相关而产生稠密整数加法群的周期，不会产生比椭圆函数更自由的非常值函数。
\ProseParagraphBreak
所谓两个周期“独立”，必须是在实数域上独立。若两个非零实周期的比为无理数，它们生成的群在实轴稠密，上一论证同样迫使常数；若比为有理数，它们只生成一维离散周期群。只有二维实独立的情形给紧的基本区域。在判定题目中的函数是否真正双周期时，应先检查这个线性代数条件，而不是数出公式中出现了两个平移数值。
\ProseParagraphBreak
格和的收敛指数二也有几何含义。一个大圆中格点数不仅是 \(O(R^2)\)，还与面积同阶：用同样大小的基本平行四边形铺平面，圆内完全包含的格胞覆盖一个稍小圆，与圆相交的格胞则落在一个稍大圆。格胞面积固定，故计数介于 \(cR^2\) 与 \(CR^2\) 之间。于是未作相消的 \(\sum|\omega|^{-2}\) 在二维格上发散，而 \(\sum|\omega|^{-3}\) 收敛。Weierstrass 级数中的减项正是跨过这个临界指数所必需的修正。
